\section*{ЗАКЛЮЧЕНИЕ}
\phantomsection
\addcontentsline{toc}{section}{Заключение}

В рамках настоящей курсовой работы выполнено численное исследование влияния регулярного дискретного микрорельефа на трибологические характеристики радиального подшипника скольжения. По итогам решены следующие задачи.

\begin{enumerate}
    \item Систематизированы теоретические представления о формировании несущего смазочного слоя и обобщены опубликованные результаты по лазерному текстурированию поверхностей трения. Показано, что определяющий вклад в повышение характеристик подшипника вносит микрогидродинамический эффект, возникающий на входных кромках микроуглублений за счёт формирования локальных сходящихся клиньев.

    \item Реализована численная схема решения уравнения Рейнольдса для подшипника конечной длины: метод конечных разностей с итерационной процедурой Гаусса - Зейделя и верхней релаксацией. Решение получено на сетке $500 \times 500$ узлов с допуском по невязке $\delta < 10^{-5}$.

    \item Получены поля давления и значения интегральных показателей~-- подъёмной силы $F$, коэффициента трения $\mu$ и объёмного расхода $Q$~-- как для гладкого подшипника, так и для подшипника с трапецеидальными (прямоугольными) углублениями (тип~10). Расчёт проведён в диапазоне относительных эксцентриситетов $\varepsilon = 0{,}05\text{ - }0{,}80$ для подшипника геометрии~A: $R = 35$~мм, $c = 50$~мкм, $L = 56$~мм.

    \item Выполнено сопоставление гладкого и текстурированного вариантов. При характерном рабочем значении $\varepsilon = 0{,}6$ зафиксированы следующие количественные изменения:
    \begin{itemize}
        \item подъёмная сила выросла на $18{,}52$\%~-- с $F = 5234{,}7$~Н до $F = 6204{,}1$~Н;
        \item коэффициент трения уменьшился на $17{,}75$\%~-- с $\mu = 0{,}01480$ до $\mu = 0{,}01217$;
        \item объёмный расход смазки увеличился на $1{,}70$\%~-- с $Q = 0{,}2356$~л/с до $Q = 0{,}2396$~л/с.
    \end{itemize}
\end{enumerate}

Полученные данные подтверждают эффективность применённой текстуры с параметрами $a = 2{,}38$~мм, $b = 2{,}18$~мм и $h_p = 10$~мкм при относительной глубине $h_p/c = 0{,}2$. Одновременное возрастание подъёмной силы и снижение коэффициента трения при пренебрежимо малом приросте расхода~-- существенный практический результат: повышение нагрузочной способности достигается без увеличения требований к маслоснабжающей системе.

Геометрия трапецеидального (прямоугольного) профиля характеризуется прямоугольной формой углубления в плане и трапецеидальным сечением: центральная зона ($r_\infty \leq 0{,}6$) имеет постоянную глубину, на периферии ($0{,}6 < r_\infty \leq 1$) глубина линейно убывает до нуля. Такая форма обеспечивает максимальную площадь зоны постоянной глубины среди рассмотренных типов, что способствует устойчивому накоплению дополнительного давления на входной по ходу вращения границе углубления и одновременно сохраняет в его центре резерв смазочной жидкости.

С учётом полученного эффекта трапецеидальный (прямоугольный) микрорельеф целесообразно рекомендовать к применению во вкладышах опор скольжения нефтегазового оборудования, эксплуатируемых в режиме гидродинамической смазки. Наибольший практический эффект ожидается в тяжелонагруженных узлах, работающих при больших относительных эксцентриситетах ($\varepsilon > 0{,}5$), где прирост подъёмной силы и снижение коэффициента трения проявляются особенно выраженно.
